\documentclass[12pt,a4paper]{article}
\usepackage[utf8]{inputenc}
\usepackage[T1]{fontenc}
\usepackage[polish]{babel}
\usepackage{amsmath, amssymb, amsthm}
\usepackage{geometry}
\usepackage{booktabs}
\usepackage{array}
\usepackage{xcolor}
\usepackage{mdframed}
\usepackage{enumitem}
\usepackage{hyperref}

\geometry{margin=2.5cm}

\definecolor{schemecolor}{RGB}{180, 60, 40}
\definecolor{boxbg}{RGB}{255, 245, 240}
\definecolor{correctgreen}{RGB}{0, 140, 60}
\definecolor{attackred}{RGB}{180, 30, 30}

\newmdenv[
backgroundcolor=boxbg,
linecolor=schemecolor,
linewidth=1.5pt,
roundcorner=5pt,
innertopmargin=10pt,
innerbottommargin=10pt,
innerleftmargin=15pt,
innerrightmargin=15pt
]{schemebox}

\newtheorem*{atak}{Atak}
\newtheorem*{wniosek}{Wniosek}

\title{
	\textbf{Lista nr 2 --- Zadanie 6}\\[0.4em]
}
\author{Kajetan Plewa}
\date{}

\begin{document}
	
	\maketitle
	\hrule
	\vspace{1em}
	
	%------------------------------------------------------------------
	\section*{Schemat podpisu}
	%------------------------------------------------------------------
	
	\begin{schemebox}
		\textbf{Parametry:} Grupa cykliczna $G = \langle g \rangle$ rzędu $q$,
		sekret $a$, klucz publiczny $A = g^a$, wiadomość $m$.
		
		\medskip
		\begin{enumerate}[label=\textbf{\arabic*.}, leftmargin=2em]
			\item \textit{Podpisujący} losuje $r$, oblicza $R = g^r$ oraz $X = g^{ar}$.
			\item Wyznacza $h = H(m, R, X)$.
			\item Oblicza $s = \bigl[h + ra\bigr] \bmod q$ i zwraca $\sigma = (R, s, X)$.
			\item \textit{Weryfikator} wyznacza $h = H(m, R, X)$ i akceptuje, gdy
			$g^s = g^h X$.
		\end{enumerate}
	\end{schemebox}
	
	%------------------------------------------------------------------
	\section{Poprawność dla uczciwego Podpisującego}
	%------------------------------------------------------------------
	
	Podstawiamy definicję $s = h + ra \pmod{q}$ do lewej strony równania
	weryfikacyjnego:
	
	\[
	g^s \;=\; g^{h + ra} \;=\; g^h \cdot g^{ra}.
	\]
	
	Ponieważ $X = g^{ar} = g^{ra}$, otrzymujemy:
	
	\[
	g^s \;=\; g^h \cdot X. \qquad {\color{correctgreen}\checkmark}
	\]
	
	\begin{wniosek}
		Schemat jest \textbf{algebraicznie poprawny} - uczciwy podpisujący zawsze
		przejdzie weryfikację.
	\end{wniosek}
	
	%------------------------------------------------------------------
	\section{Czy dodanie $X$ do hasha naprawia problem z zadania~5?}
	%------------------------------------------------------------------
	
	W zadaniu~5 (gdzie $h = H(m, R)$) atakujący wykonywał:
	\[
	\text{wybierz } s, \quad \text{oblicz } h = H(m, R), \quad
	\text{ustaw } X = g^{s-h}.
	\]
	
	W obecnym schemacie $h$ zależy od $X$, a atakujący chciałby ustawić
	$X = g^{s-h}$, co tworzy \textbf{zależność kołową}:
	\[
	h = H(m, R, X) = H\!\left(m, R,\, g^{s - h}\right),
	\]
	czyli $h$ pojawia się po obu stronach - nie można go wyznaczyć
	tym podejściem.
	
	\medskip
	Zatem dodanie $X$ do wejścia funkcji skrótu \emph{blokuje konkretny atak
		z zadania~5}. 
	
	%------------------------------------------------------------------
	\section{Atak fałszerstwa bez znajomości $a$}
	%------------------------------------------------------------------
	
	Kluczowa obserwacja: równanie weryfikacyjne $g^s = g^h X$ \textbf{nigdzie
		nie zawiera klucza publicznego $A = g^a$}. Weryfikator nigdy nie sprawdza,
	czy $X = A^r$, co pozwala na następujący atak.
	
	\begin{atak}
		Atakujący chce sfałszować podpis pod dowolną wiadomością $m$:
		\begin{enumerate}[label=\arabic*.]
			\item Wybierz dowolne $r' \in \mathbb{Z}_q$, ustaw $R = g^{r'}$.
			\item Wybierz dowolne $k \in \mathbb{Z}_q$, ustaw $X = g^k$.
			\item Oblicz $h = H(m, R, X)$.
			\item Ustaw $s = h + k \pmod{q}$.
			\item Zwróć $\sigma = (R,\, s,\, X)$.
		\end{enumerate}
	\end{atak}
	
	\noindent\textbf{Weryfikacja:}
	\[
	g^s = g^{h+k} = g^h \cdot g^k = g^h \cdot X. \qquad {\color{correctgreen}\checkmark}
	\]
	Atak jest więc skuteczny dla każdej wiadomości.
	
	%------------------------------------------------------------------
	\section{Porównanie z zadaniem~5}
	%------------------------------------------------------------------
	
	\begin{center}
		\renewcommand{\arraystretch}{1.4}
		\begin{tabular}{lcc}
			\toprule
			\textbf{Właściwość} & \textbf{Zadanie 5} & \textbf{Zadanie 6} \\
			\midrule
			Hash & $H(m, R)$ & $H(m, R, X)$ \\
			Atak z zad.\ 5 działa? & {\color{correctgreen}tak} & {\color{attackred}nie --- kołowość} \\
			Istnieje atak  & {\color{attackred}tak} & {\color{attackred}tak (wybór $k$)} \\
			$A$ w równaniu weryfikacyjnym & {\color{attackred}nie} & {\color{attackred}nie} \\
			Poziom bezpieczeństwa & 0 bitów & 0 bitów \\
			\bottomrule
		\end{tabular}
	\end{center}
	
	\medskip
	Zmiana $H(m,R) \to H(m,R,X)$ jest \textbf{pozorna}: uszczelnia jeden
	konkretny wektor ataku.  Żeby faktycznie rozwiązać problem tego schematu należy dodać weryfikacje $X$. Ponieważ obecnie weryfikator wierzy bezgranicznie inputowi podpisującego, sekret jest niewykorzystywany a schemat niebezpieczny.
\end{document}